\documentclass[12pt,letterpaper]{article}
\usepackage[margin=0.8in]{geometry}
\usepackage{amsmath,amssymb}
\pagestyle{empty}
\setlength{\parindent}{0pt}
\setlength{\parskip}{0.7em}
% practice set 2: LQ3
% Q4--8 retained from Practice mock - 2026-10-18.
% Q9: Generated LQ3-SET2-PE-01.
% Q10: Course Quantum practice problems, Compton scattering 2(b).
% Q11: Generated LQ3-SET2-LIGHT-01.
\begin{document}

4. Two bullets collide head-on. The first bullet is going with speed $0.1c$ and has rest mass $5\,\mathrm{g}$. The second bullet is going with speed $0.2c$ and has rest mass $3\,\mathrm{g}$. Assume that after the collision, the two bullets form one stuck-together mass.

What are the rest mass and velocity of the resulting object? State your positive direction.

\vfill
\newpage

5. A star emits light with wavelength $500\,\mathrm{nm}$ in its own rest frame. Its speed relative to Earth is $v$. Give your answers in terms of $v$ and $c$.

(a) If the star is moving directly away from Earth, what wavelength does an observer on Earth detect?

\vfill

(b) If the star is moving perpendicular to the line of sight instead, what is the observed wavelength? The line of sight is defined in the Earth frame at emission.

\vfill
\newpage

6. A particle of rest mass $10m$ moves at $+0.60c$ along the laboratory $x$-axis. It decays into two particles, each of rest mass $3m$. In the parent's rest frame, daughter A moves along $+x$ and daughter B along $-x$. No other particles are produced.

Express energies in multiples of $mc^2$, momenta in multiples of $mc$, and velocities in multiples of $c$.

(a) Find the total energy and signed $x$-momentum of each daughter in the parent rest frame.

\vfill

(b) Find the total energy and signed $x$-momentum of each daughter in the laboratory frame.

\vfill

(c) Find each daughter's signed laboratory velocity.

\vfill
\newpage

7. Treat the electrons as free and initially stationary. You may use $h/(m_ec)=0.00243\,\mathrm{nm}$.

(a) X-rays of wavelength $0.100\,\mathrm{nm}$ scatter directly backwards off electrons in a thin sheet of metal. What is the wavelength of the scattered photons?

\vfill

(b) An optical photon of wavelength $500\,\mathrm{nm}$ scatters off an electron at an angle of $45^\circ$. How big is the wavelength shift as a fraction of the original wavelength?

\vfill
\newpage

8. Two proposed experiments each begin with two particles of rest mass $m$ and total laboratory energy $6mc^2$.

In experiment A, the particles move towards one another with equal and opposite momenta.

In experiment B, one particle is initially stationary.

In each experiment, the particles combine into a single product, with no radiation or other products.

Which experiment produces the larger product rest mass? Calculate the product rest mass in each case and briefly explain the difference.

\vfill
\newpage

9. A metal has work function $2.20\,\mathrm{eV}$. Monochromatic light of wavelength $400\,\mathrm{nm}$ falls on its surface. Use $hc=1240\,\mathrm{eV\,nm}$.

(a) Find the maximum kinetic energy of the emitted electrons and the magnitude of the stopping potential.

\vfill

(b) Find the longest wavelength that can produce photoelectric emission.

\vfill

(c) The intensity is doubled while the wavelength remains $400\,\mathrm{nm}$. What happens to the maximum electron kinetic energy and the number of electrons emitted per second? Assume a constant emission efficiency and no saturation.

\vfill
\newpage

10. A photon scatters off a free electron initially at rest. You may use $h/(m_ec)=0.00243\,\mathrm{nm}$.

(a) At what scattering angle is the change in the photon's wavelength greatest?

\vfill

(b) In a Compton scattering experiment, some photons are observed to increase their wavelength by $1\%$. What is the maximum possible wavelength of the incoming photons?

\vfill
\newpage

11. A uniform monochromatic laser beam in vacuum has wavelength $600\,\mathrm{nm}$, power $3.00\,\mathrm{mW}$, and cross-sectional area $2.00\,\mathrm{mm}^2$. Use $h=6.626\times10^{-34}\,\mathrm{J\,s}$ and $c=3.00\times10^8\,\mathrm{m/s}$.

(a) Find its frequency, intensity, and the number of photons crossing the whole beam cross-section each second.

\vfill

(b) Its electric-field amplitude is doubled while its frequency and cross-sectional area remain unchanged. By what factors do its intensity and photon arrival rate change?

\vfill

(c) In a separate change from the original beam, the wavelength is halved while the power and cross-sectional area remain unchanged. By what factors do the photon energy and photon arrival rate change?

\vfill
\end{document}
